% kernel/engineering_decisions.tex

\chapter{From Evaluated Proposal to Authorized State}
\label{chap:engineering-decisions}

The Berlinish composition has survived the required geometric, realization,
speed, cant, and evidence checks. It is now an applicable candidate for the
stated purpose. It is still not the accepted state of the alignment.

\section{Proposal, candidate, and evaluation}

\begin{definition}[Proposal]
\label{def:proposal}
A proposal is a candidate engineering change submitted for evaluation before
acceptance or rejection.
\end{definition}

\begin{definition}[Candidate solution]
\label{def:candidate-solution}
A candidate solution is a proposed engineering state or configuration that
may be evaluated against applicable knowledge, constraints, and preferences.
\end{definition}

axtranNew searches or solves within the expressible transition space. The
engineering object model carries the resulting proposal and its relation to
the alignment. Evaluation determines purpose-relative sufficiency. These are
separate responsibilities: successful generation does not imply admissibility,
and admissibility does not select among all admissible alternatives.

\section{Authority changes engineering status}

The active Kernel fixes the boundary that evaluation does not exercise
decision authority; it does not yet establish a complete canonical identity
for every Engineering Decision. For this explanatory sequence, the following
narrow operational formulation is sufficient:

\begin{definition}[Engineering decision]
\label{def:engineering-decision}
An Engineering Decision is an authorized acceptance, rejection, or selection
that resolves an engineering proposal or alternative within explicitly stated
qualifications.
\end{definition}

The entitled authority may accept the candidate as a successor state, reject
it, request further evidence, or select another candidate. The decision names
its scope, subject, evidence, applicable knowledge, responsible authority, and
effect. It does not erase the predecessor, rejected candidates, survey, or
calculation record.

If accepted, the new state remains related to the same alignment when the
authorized identity judgment preserves object continuity. If the change is
governed as a replacement object, that relation must be stated instead. The
Kernel does not supply a universal geometric threshold between these outcomes;
the solver therefore cannot manufacture one.

\section{Return to the transition}

The running example has crossed the complete meaning chain:

\begin{enumerate}
\item Berlinish expressed an ordered transition construction.
\item axtranNew or an equivalent service calculated a parameterized proposal.
\item the proposal was related to an identifiable alignment and candidate
  state while preserving its constructive composition;
\item Metric Realization and Realization Context made it spatially measurable;
\item representations communicated it and provenance separated calculation,
  import, derivation, and observation;
\item applicable knowledge and purpose-relative evaluation tested its
  sufficiency; and
\item an entitled authority, if it chose to do so, admitted the successor
  state.
\end{enumerate}

At no point did computation create authority, representation acquire identity,
observation become physical reality, or evaluation constitute decision. This
is the engineering meaning carried forward into Geometry and Curvature: the
mathematics will now operate on identifiable, qualified states rather than on
anonymous curves.
